\section{远场光子是总损失中的哪一部分}
\label{chap:feqo-cl}

平面界面的电子可激发很强的近场，却因平移动量过大而没有远场光。若样品有限、边缘破坏平移对称性，或另有光栅提供动量，式~\eqref{eq:feqo-cl-collection} 才可能给非零计数。实际计算先用同一电流源求 \(\mathbf E=\ii\mu_0\omega\mathbf G\mathbf j_e\)，再取其远场系数 \(\mathbf F(\hat{\mathbf n},\omega)\)。这个顺序保留了材料响应与观测投影的区别。

把这个操作写成核的形式，就能看见探测器究竟投影了什么。对远处的观测点定义 \(\mathbf G_\infty(\hat{\mathbf n},\mathbf r';\omega)=\lim_{r\to\infty}r e^{-\ii kr}\mathbf G(r\hat{\mathbf n},\mathbf r';\omega)\)，前提是该方向存在通常的 \(e^{\ii kr}/r\) 远场展开。则
\[
\mathbf F(\hat{\mathbf n},\omega)
=\ii\mu_0\omega\int\dd^3r'\,
\mathbf G_\infty(\hat{\mathbf n},\mathbf r';\omega)
\mathbf j_e(\mathbf r',\omega).
\]
同一电子电流在式~\eqref{eq:feqo-trajectory-green-kernel} 中与 \(\operatorname{Im}\mathbf G\) 作二次型，给所有损失通道；在这里先由 \(\mathbf G_\infty\) 投影到指定方向，再取模平方，只给能到远场的部分。吸收模与导波模没有该方向的远场系数，因而不会凭空出现在 CL 计数中。

对包住样品的闭合曲面，时域 Poynting 定理把电子交出的能量分成介质吸收、向外辐射和最终仍储存在局域场中的三项。若脉冲前后局域储能回到原值，频率分辨的能量收支可写为 \(W_{\rm loss}=W_{\rm abs}+W_{\rm rad}\)，但只有在同一物理区域、同一入射条件且模式间交叉项正确处理后才可逐频使用。定义 \(\eta_{\rm rad}=W_{\rm rad}/(W_{\rm rad}+W_{\rm abs})\) 与 \(\eta_{\rm coll}=W_{\rm coll}/W_{\rm rad}\)，二者分别回答“多少能量成为远场光”和“多少远场光进入本探测器”。光学数值孔径只能改变第二项，不能恢复已变成热的能量。

由式~\eqref{eq:feqo-cl-collection}，改变收集角不一定只改变总强度：若两个模式有不同角分布，探测器对它们的相对权重也会改变。比较电子损失谱与发光谱时，必须给出收集角、偏振与效率；把一个未标定的发光峰直接解释成材料态密度会混入仪器选择。

一个小颗粒的偶极近似可把角度权重算完。假设颗粒远小于波长，电子在颗粒位置产生的频域驱动场为 \(\mathbf E_e(\omega)\)，颗粒的线性极化率为体积量纲的复数 \(\alpha(\omega)\)。若响应只沿 \(z\) 方向，诱导偶极矩为 \(p_z(\omega)=\varepsilon_0\alpha(\omega)E_{e,z}(\omega)\)。真空偶极远场满足
\[
\mathbf F(\hat{\mathbf n},\omega)
=\frac{k^2}{4\pi\varepsilon_0}
\bigl[\hat{\mathbf n}\times(\mathbf p\times\hat{\mathbf n})\bigr],
\qquad k=\omega/c.
\]
把 \(\mathbf p=p_z\hat{\mathbf z}\) 代入，\(|\mathbf F|^2=k^4|p_z|^2\sin^2\theta/(16\pi^2\varepsilon_0^2)\)。式~\eqref{eq:feqo-cl-collection} 因而预言角分布为 \(\sin^2\theta\)，而不是各方向相同。若理想探测器收集以 \(+z\) 为轴、半角 \(\theta_{\max}\) 的圆锥，且效率在圆锥内为一，则因 \(\dd\Omega=2\pi\sin\theta\dd\theta\)，收集到的辐射比例为
\begin{equation}
\eta_{\rm cone}=
\frac{\int_0^{\theta_{\max}}\sin^3\theta\dd\theta}
{\int_0^\pi\sin^3\theta\dd\theta}
=\frac{2-3\cos\theta_{\max}+\cos^3\theta_{\max}}4.
\label{eq:feqo-dipole-collection-fraction}
\end{equation}
半球收集给 \(1/2\)，很窄的轴向圆锥因偶极轴上恰有辐射零点而远小于它的立体角比例。这个算例从已知电子驱动场 \(E_{e,z}\) 经材料极化率算到远场角分布与收集计数；它依赖偶极、线性响应和真空远场条件，不能直接用在尺寸可比波长的颗粒或复杂衬底上。

\begin{exercise}
设一个模式把总能量的 \(80\%\) 送到远场，探测器收集其中 \(25\%\)。求 \(\eta_{\rm rad}\eta_{\rm coll}\)，并说明改变数值孔径能否改变模式内部的吸收功率。
\end{exercise}
